\section{Datasets and roles}

The benchmark uses each dataset for a specific scientific role rather than
treating all resources as interchangeable evidence for a single claim.

\paragraph{Synthetic known-truth cases.}
Four synthetic cases validate the Mechanistic Identifiability Score under
controlled truth: directed truth, common drive, structure without temporal
direction, and timing without region-specific topology.

\paragraph{Allen Visual Behavior Neuropixels.}
Allen VBN is used as a negative mechanistic-identifiability control. The current
analysis includes 20 successful sessions in the confirmation phase, 20 sessions
and 20 unique mice in the anatomical-connectivity test, and 21 sessions with 19
unique mice in the directed-identifiability test.

\paragraph{Sensorium static.}
Sensorium 2022 static is used as an auxiliary predictive and reliability case.
The local package includes seven validation mice and five repeated-test mice.
Across the repeated-test mice, median reliability is 0.642 and median best
predictive correlation is 0.341.

\paragraph{Dynamic Sensorium.}
Dynamic Sensorium is used as a temporal predictive stress test. The current
package includes five current public mice and five legacy OOD mice. Transparent
baselines include mean response, summary features, temporal filterbank,
temporal SVD, random features, a compact PyTorch MLP, and a bounded official
Sensorium model family. The official bounded baseline is integrated and
evaluated but remains non-SOTA and non-Q1-qualified.

\paragraph{MICRONS CAVE.}
MICRONS is the main empirical evidence layer. The Q1 package contains three
non-overlapping CAVE windows: a discovery cohort and two offset hold-outs,
together covering 2991 co-registered units, 6575 synapses, and 5943 connected
directed pairs.
